\chapter{The Circumgalactic Medium, the Intergalactic Medium, and the Cosmic Web}
\label{chap:cgm-igm}

\section{A working map of diffuse baryons}

The gaseous Universe is not naturally divided into boxes. Gas crosses the
interstellar medium (ISM), circumgalactic medium (CGM), intragroup medium, and
intergalactic medium (IGM); the same parcel may be accreted, form stars, be
ejected, cool, mix, and return. The names are nevertheless useful if they are
treated as operational definitions rather than immutable phases. In this
chapter, the \emph{CGM} is gas outside the bright stellar body and cold
star-forming disk but inside the region dynamically associated with a galaxy's
dark-matter halo. The \emph{IGM} is gas outside individual halos. The
\emph{cosmic web} is the evolving network of sheets, filaments, nodes, and
voids created by gravitational instability. These definitions deliberately
leave a transition zone: around groups, interacting systems, and splashback
regions, assignment to one galaxy can be physically ambiguous
\citep{Tumlinson2017,FaucherGiguereOh2023}.

A convenient halo scale is defined by
\begin{equation}
 M_{\Delta}=\frac{4\pi}{3}\,\Delta\,\rho_{\rm ref}(z)R_{\Delta}^{3},
 \label{eq:halo-definition}
\end{equation}
where the reference density may be the critical or mean matter density.
Authors who quote $R_{200\rm c}$, $R_{200\rm m}$, or a "virial radius" are
not necessarily using the same boundary. The conversion can be tens of
percent and should never be hidden when normalizing impact parameters.
Associated velocity and temperature scales are
\begin{align}
 V_{\Delta} &= \left(\frac{GM_{\Delta}}{R_{\Delta}}\right)^{1/2},\\
 T_{\rm vir} &\simeq \frac{\mu m_{\rm p}V_{\Delta}^{2}}{2k_{\rm B}},
 \label{eq:tvir}
\end{align}
with mean molecular weight $\mu$. For an ionized, Milky-Way-mass halo,
$T_{\rm vir}$ is of order $10^{6}$ K. This number is a guide, not a claim
that the halo is isothermal.

The most important conceptual correction to a one-phase cartoon is that
temperature labels do not define unique structures. "Cool" gas near
$10^{4}$ K may occur in infalling filaments, stripped satellite wakes,
radiatively cooled wind material, turbulent mixing layers, or self-shielded
clouds. "Warm" gas near $10^{5}$--$10^{6}$ K may be cooling through a
short-lived temperature interval rather than occupying a stable volume.
"Hot" gas can be supported approximately hydrostatically, driven by winds,
or shock heated during accretion. Observed velocity coincidence between two
ions does not prove that they are co-spatial. A spectrograph collapses a
three-dimensional, time-dependent medium into optical depth versus one
velocity coordinate.

\section{Why the CGM is central to galaxy evolution}

A galaxy's stellar mass is only the integrated residue of a larger baryon
cycle. Cosmological inflow supplies gas; cooling and angular-momentum
transport deliver some of it to the ISM; stars and black holes return mass,
metals, momentum, cosmic rays, and radiation; satellites donate gas by winds
and stripping; and some material escapes or is exchanged between halos.
Consequently, the CGM is simultaneously a fuel reservoir, a feedback
calorimeter, and a record of environmental processing
\citep{Tumlinson2017,FaucherGiguereOh2023,Hafen2019}.

For a control surface at radius $r$, a formal instantaneous mass flux is
\begin{equation}
 \dot{M}(r)=\int \rho(\boldsymbol{x})
        \left[\boldsymbol{v}(\boldsymbol{x})-\boldsymbol{v}_{\rm frame}\right]
        \cdot \hat{\boldsymbol{r}}\,r^{2}\,{\rm d}\Omega .
 \label{eq:mass-flux}
\end{equation}
This equation is straightforward in a simulation and extremely difficult to
infer from a single sightline. An absorber supplies a column density and a
line-of-sight velocity, but generally not its distance along the sightline,
solid angle, transverse velocity, or time history. Thus labels such as
"inflow" and "outflow" should be probabilistic unless geometry supplies
additional evidence. Major-axis/minor-axis trends, metallicity, comparison
with the escape speed, resolved emission, or multiple nearby sightlines can
strengthen the interpretation, but none is individually decisive.

The cool-CGM mass inferred from absorption illustrates the inference chain.
For a hydrogen column profile $N_{\rm H}(R)$ and azimuthal covering fraction
$f_{\rm c}(R)$,
\begin{equation}
 M_{\rm cool}(<R_{\rm max}) =
 2\pi \mu_{\rm H}m_{\rm p}
 \int_{0}^{R_{\rm max}} N_{\rm H}(R) f_{\rm c}(R)R\,{\rm d}R .
 \label{eq:cgm-mass}
\end{equation}
The directly constrained quantity is often $N_{\rm H\,I}$, not $N_{\rm H}$.
Writing $f_{\rm H\,I}=N_{\rm H\,I}/N_{\rm H}$ gives
$N_{\rm H}=N_{\rm H\,I}/f_{\rm H\,I}$. The neutral fraction can be tiny and
depends on density, spectrum and amplitude of the ionizing field, temperature,
shielding, and non-equilibrium history. Large masses inferred for cool halo
gas are therefore important but model conditioned
\citep{Werk2014,Prochaska2017}. A robust paper reports which part of
Equation~\ref{eq:cgm-mass} is measured and which part is assumed.

\section{Thermal competition and multiphase structure}

Two timescales provide a useful first organization. An approximate radiative
cooling time is
\begin{equation}
 t_{\rm cool} \simeq
 \frac{\tfrac{3}{2}n_{\rm tot}k_{\rm B}T}
      {n_{\rm e}n_{\rm H}\Lambda(T,Z,J_{\nu})},
 \label{eq:cooling-time}
\end{equation}
where $\Lambda$ includes radiative losses and, if desired, photoheating. A
local free-fall estimate is
\begin{equation}
 t_{\rm ff}(r)=\left[\frac{2r}{g(r)}\right]^{1/2}.
 \label{eq:freefall}
\end{equation}
When shock-heated gas cools slowly compared with its dynamical time, an
extended hot atmosphere can persist. When cooling is rapid, cold streams,
condensation, and radiative interfaces become important. The ratio
$t_{\rm cool}/t_{\rm ff}$ is a diagnostic, not a universal one-number
criterion; rotation, turbulence, uplift, conduction, magnetic topology,
cosmic-ray pressure, and ongoing forcing all alter the outcome
\citep{FaucherGiguereOh2023}.

Pressure balance is similarly informative but not guaranteed:
\begin{equation}
 P_{\rm tot}=P_{\rm th}+P_{\rm turb}+P_{\rm B}+P_{\rm CR}+\cdots .
\end{equation}
A cool component can have $n_{\rm cool}T_{\rm cool}\neq
n_{\rm hot}T_{\rm hot}$ while remaining confined by non-thermal pressure,
ram pressure, gravity, or time dependence. An apparent pressure discrepancy
between photoionization models and a coronal model may therefore diagnose
missing physics, but it may also expose an incorrect ionizing background or
an assumption that non-coincident velocity components form one cloud.

Interfaces are disproportionately visible. Conductive boundaries,
radiative shocks, and turbulent mixing can generate ions such as
$\mathrm{C\,IV}$, $\mathrm{N\,V}$, $\mathrm{O\,VI}$, and
$\mathrm{Ne\,VIII}$ while containing little of the total mass. In
collisional ionization equilibrium, each ion fraction has a characteristic
temperature range, but gas that cools faster than it recombines retains an
ionization memory. Photoionization can also produce high ions at much
lower kinetic temperature if the density is sufficiently small. Ion names
are therefore not thermometers by themselves
\citep{GnatSternberg2007,SutherlandDopita1993}.

\section{A tracer ladder for circumgalactic gas}

\subsection{Neutral hydrogen}

$\mathrm{H\,I}$ is usually the entry point because the Lyman series spans an
enormous range of column density. Ly$\alpha$ is sensitive but saturates
readily. Higher-order Lyman lines have smaller oscillator strengths and
recover information hidden in saturated Ly$\alpha$. At sufficiently high
$N_{\rm H\,I}$, the continuum optical depth at the Lyman edge,
\begin{equation}
 \tau_{912}=N_{\rm H\,I}\sigma_{912},
\end{equation}
constrains the neutral column even when line cores are black. At still larger
columns, Lorentzian damping wings determine $N_{\rm H\,I}$. The lesson is to
design coverage of a \emph{ladder}, not a single line.

Absorber names encode approximate column regimes, but boundaries vary with
author and purpose. The Ly$\alpha$ forest generally denotes mostly
photoionized systems below the classical Lyman-limit threshold. Partial
Lyman-limit systems and Lyman-limit systems occupy the transition to optical
thickness near 1 Ryd. Sub-damped and damped Ly$\alpha$ systems contain
progressively larger neutral columns; the conventional DLA boundary is
$N_{\rm H\,I}=2\times10^{20}\ {\rm cm^{-2}}$
\citep{Wolfe2005}. These are observational classes, not one-to-one maps onto
IGM, CGM, and ISM.

\subsection{Low and intermediate ions}

$\mathrm{Mg\,II}$, $\mathrm{C\,II}$, $\mathrm{Si\,II}$,
$\mathrm{Fe\,II}$, and $\mathrm{O\,I}$ preferentially trace cool or
low-ionization material. Their diagnostic value differs:
\begin{itemize}
  \item $\mathrm{O\,I}$ is coupled to $\mathrm{H\,I}$ by charge exchange over
  useful conditions, making O/H comparatively direct in neutral gas.
  \item $\mathrm{Si\,II}$ and $\mathrm{Fe\,II}$ have several transitions
  spanning oscillator strengths, valuable for diagnosing saturation.
  \item Refractory Fe and Si can be depleted onto dust, so gas-phase abundance
  ratios need not equal nucleosynthetic ratios.
  \item Mg II is observable from the ground after modest redshifting and has
  built large statistical samples, but equivalent width mixes column,
  velocity spread, covering fraction, and saturation.
\end{itemize}

Intermediate ions such as $\mathrm{Si\,III}$ and $\mathrm{C\,III}$ are often
strong because they can be abundant in photoionized gas, but their strongest
transitions saturate easily. A precise equivalent width is not necessarily a
precise column density.

\subsection{High ions and hot gas}

$\mathrm{C\,IV}$, $\mathrm{N\,V}$, $\mathrm{O\,VI}$, and
$\mathrm{Ne\,VIII}$ probe harder ionization conditions and often reveal
interfaces or warm-hot gas. O VI is especially important because its
$1031.93$ and $1037.62$ \AA\ doublet is accessible to far-UV spectroscopy at
low redshift and to optical spectroscopy at sufficiently high redshift.
Strong, extended O VI around star-forming galaxies established a major metal
reservoir \citep{Tumlinson2011}. Yet converting $N_{\rm O\,VI}$ into oxygen
or total gas mass requires both the ion fraction $f_{\rm O\,VI}$ and
metallicity:
\begin{equation}
 N_{\rm H} =
 \frac{N_{\rm O\,VI}}
 {f_{\rm O\,VI}\,({\rm O/H})_{\odot}\,10^{[{\rm O/H}]}} .
 \label{eq:ovi-correction}
\end{equation}
A factor-of-ten uncertainty in metallicity becomes a factor-of-ten
uncertainty in mass.

At $T\gtrsim10^{6}$ K, abundant ions emit and absorb soft X-rays. Oxygen
$\mathrm{O\,VII}$ and $\mathrm{O\,VIII}$ are central WHIM and halo tracers;
the Fe-L complex diagnoses hotter, metal-rich plasma. Emission scales
approximately as $\int n_{\rm e}n_{\rm H}\Lambda_E\,{\rm d}\ell$ and is
weighted toward dense clumps. Absorption scales as $\int n_i\,{\rm d}\ell$
and is more sensitive to diffuse columns but requires an exceptionally bright
background source and exquisite control of instrumental structure. The two
methods do not sample the same density weighting.

\section{The IGM as a fluctuating medium}

\subsection{Photoionized forest}

The classical cloud picture has largely been replaced by a continuous-density
view. Much Ly$\alpha$ forest absorption arises in mildly overdense,
photoionized sheets and filaments whose peculiar velocities, temperature,
thermal broadening, and Hubble expansion shape the spectrum
\citep{Rauch1998,McQuinn2016}. For highly ionized hydrogen in equilibrium,
\begin{equation}
 n_{\rm H\,I}\Gamma_{\rm H\,I}
 \simeq n_{\rm e}n_{\rm p}\alpha_{\rm A}(T),
 \label{eq:photoequilibrium}
\end{equation}
so the neutral fraction is extremely sensitive to the photoionization rate
$\Gamma_{\rm H\,I}$ and moderately sensitive to temperature through the
recombination coefficient. The transmitted flux is
\begin{equation}
 F(v)=C(v)\exp[-\tau(v)]+N(v),
 \label{eq:forest-flux}
\end{equation}
with continuum $C$, absorption field $\tau$, and noise $N$. In forest
cosmology, continuum fitting is not merely cosmetic: fitting away
long-wavelength modes distorts the measured correlation function and must be
represented in the forward model.

A useful analytic argument sets an absorber's characteristic scale near a
local Jeans or filtering scale. Combining photoionization equilibrium with
that scale produces a relation between $N_{\rm H\,I}$ and gas overdensity
\citep{Schaye2001}. The relation explains why the column-density distribution
can encode the matter distribution, but its coefficients depend on
temperature, ionizing background, gas fraction, and dynamical assumptions.
It should be used as a calibrated scaling, not an exact cloud ruler.

\subsection{Mean opacity, thermal history, and reionization}

Even a small neutral fraction produces substantial Ly$\alpha$ opacity because
the transition is strong. The Gunn--Peterson argument therefore constrains
diffuse neutral hydrogen far below the fraction needed for a predominantly
neutral Universe \citep{GunnPeterson1965}. A dark trough at high redshift does
not by itself measure a unique neutral fraction: once Ly$\alpha$ saturates,
small changes in transmission can correspond to a broad range of neutral
fractions. Higher Lyman lines, damping-wing shape, proximity-zone structure,
and ensembles of sightlines are needed.

The forest also remembers heat input. Photoheating during hydrogen and
helium reionization establishes a temperature--density relation often written
\begin{equation}
 T(\Delta)=T_0\Delta^{\gamma-1},
\end{equation}
for a restricted low-density regime, where
$\Delta=\rho_{\rm b}/\bar{\rho}_{\rm b}$. Thermal broadening, pressure
smoothing, and peculiar velocities affect line widths and small-scale flux
power differently. Inferring $T_0$ and $\gamma$ therefore requires synthetic
spectra with the instrumental line-spread function and the same continuum and
masking operations as the data.

\subsection{The warm-hot IGM}

As large-scale structure grows, accretion and merger shocks heat some diffuse
gas to $10^{5}$--$10^{7}$ K, the conventionally defined warm-hot
intergalactic medium (WHIM) \citep{CenOstriker1999,Dave2001}. No single tracer
is complete:
\begin{description}
  \item[Broad Ly$\alpha$ absorbers] are metallicity independent in principle,
  but blends, continuum errors, and non-thermal broadening can mimic large
  Doppler parameters.
  \item[O VI and Ne VIII] are powerful in the UV, but their baryon conversion
  depends on metallicity and ionization history.
  \item[O VII and O VIII] access hotter gas in X-rays, but line equivalent
  widths are tiny and claimed detections demand independent confirmation,
  careful trials accounting, and stable calibration.
  \item[Thermal Sunyaev--Zel'dovich and dispersion measures] respond to
  electrons rather than a metal abundance, but projection and halo separation
  are nontrivial.
\end{description}
Historically, "missing baryons" has meant baryons permitted by early-Universe
constraints but absent from low-redshift inventories, not non-baryonic dark
matter. Census estimates evolve with surveys and ionization corrections
\citep{Bregman2007,Shull2012,PerouxHowk2020}. Individual X-ray WHIM claims,
including influential O VII results, should be presented with their stated
significance and subsequent replication status rather than as closure of the
census \citep{Nicastro2018}.

\section{Metals as history, not merely abundance}

Metals in diffuse gas prove prior processing by stars, but translating them
into a transport history is difficult. The logarithmic abundance is
\begin{equation}
 [{\rm X/H}]=
 \log_{10}\left(\frac{N_{\rm X}}{N_{\rm H}}\right)
 -\log_{10}\left(\frac{\rm X}{\rm H}\right)_{\odot},
 \label{eq:metallicity}
\end{equation}
after ionization and depletion corrections. The adopted solar reference
must be stated \citep{Asplund2021}. A metallicity posterior inherits
uncertainty from $N_{\rm H\,I}$, the radiation field, gas density, relative
abundance pattern, dust depletion, phase assignment, and atomic data.

Low metallicity can support an accretion interpretation, while enhanced
$\alpha$/Fe or N/O can encode enrichment timescales. However, metal-poor gas
may be old IGM material mixed only weakly with winds; metal-rich gas may arise
from outflow, recycling, stripping, or enrichment by another galaxy.
Metallicity alone does not provide the sign of radial velocity. Observed
metallicity distributions are also selection functions: metal-selected
samples and H I-selected samples answer different questions. Early evidence
for a bimodal low-redshift Lyman-limit metallicity distribution was
scientifically influential \citep{Lehner2013}; later analyses should be
compared at matched H I range, redshift, sensitivity, and priors rather than
reduced to a slogan.

\section{From pencil beams to maps}

Absorption is extraordinarily sensitive because detectability does not suffer
the same density-squared penalty as emission, but it sparsely samples the sky.
Emission maps morphology but is surface-brightness limited and vulnerable to
sky subtraction, scattered light, and resonant transfer. The strongest
programs combine:
\begin{enumerate}
  \item a background quasar or galaxy grid for column density and kinematics;
  \item deep galaxy redshifts to establish environment and impact parameters;
  \item integral-field emission mapping for morphology and velocity;
  \item X-ray, millimeter, or radio measurements for phases missed in the UV;
  \item synthetic observations of simulations passed through the same
  selection, resolution, and noise.
\end{enumerate}

At high redshift, extended Ly$\alpha$ emission can illuminate cool gas around
galaxies and quasars. Ly$\alpha$ is resonant: a photon can emerge far from
where it was produced and at a frequency altered by scattering through a
moving medium. Morphology and line centroid therefore cannot be read as a
direct emissivity or velocity map without radiative-transfer modeling.
Non-resonant recombination lines and metal lines provide critical controls.

Dense background-source fields enable tomography. Each sightline remains a
one-dimensional convolution of density, ionization, peculiar velocity, and
thermal state, but regularized reconstruction can reveal coherent filaments
and voids. On much larger scales, surveys treat transmitted forest flux as a
biased tracer and measure its auto-correlation and cross-correlation with
quasars. DESI has demonstrated the statistical power of hundreds of
thousands of forest spectra for baryon acoustic oscillations
\citep{DESI2025LyA}. This regime is conceptually continuous with single-line
analysis: both require a forward model for the continuum, absorption,
instrument, masking, and selection.

\section{A disciplined interpretation protocol}

For any diffuse-gas claim, write the chain in the following order.
\begin{enumerate}
  \item \textbf{Observed:} normalized flux, surface brightness, equivalent
  width, profile shape, angular separation, and their covariance.
  \item \textbf{Instrumental:} spectral and spatial line-spread functions,
  wavelength zero point, sensitivity, background model, detector artifacts,
  aperture, and completeness.
  \item \textbf{Atomic inference:} transition identity, oscillator strength,
  column density, excitation, saturation assessment, and blend model.
  \item \textbf{Plasma inference:} ionization mechanism, density,
  temperature, radiation field, abundance pattern, depletion, and phase
  association.
  \item \textbf{Geometric inference:} distance, path length, covering
  fraction, opening angle, inclination, and association with a galaxy or
  filament.
  \item \textbf{Derived physics:} total mass, flux, cooling rate, metal mass,
  energy, and comparison with binding or feedback budgets.
\end{enumerate}
Uncertainty normally grows down this list. Reporting a precise equivalent
width with a loosely constrained total mass is not contradictory; it is an
honest consequence of the inference hierarchy.

\section{Common category errors}

\begin{itemize}
  \item \textbf{Ion equals phase.} O VI can arise from collisionally ionized,
  photoionized, non-equilibrium cooling, or mixed gas.
  \item \textbf{Velocity equals radial motion.} A blueshift is not
  automatically an outflow if systemic redshift, inclination, rotation, and
  environment are uncertain.
  \item \textbf{Equivalent width equals column.} Saturated equivalent width
  can increase through velocity spread with little sensitivity to column.
  \item \textbf{Emission and absorption sample the same gas.} Their density
  weighting and selection functions differ.
  \item \textbf{A halo boundary is observed.} $R_{\rm vir}$ is convention and
  model; gas dynamics need not change abruptly there.
  \item \textbf{Metallicity is origin.} Enrichment constrains history but
  rarely identifies a unique transport channel.
  \item \textbf{A simulation cell is a cloud.} Numerical resolution, mixing,
  and subgrid feedback can alter synthetic absorber structure; convergence of
  galaxy properties does not guarantee convergence of CGM interfaces.
\end{itemize}

\section{Research frontiers}

Several measurements would transform the field. Closely spaced sightlines
can determine transverse coherence and cloud sizes without assuming a path
length. Joint UV and X-ray spectroscopy can connect cool, interface, and
virial-temperature gas. Deep, absolutely calibrated emission mapping can
measure faint cooling and fluorescence. Time-domain absorption against
moving or variable background sources can probe small-scale structure.
Finally, hierarchical analyses that fit populations rather than one absorber
at a time can propagate non-detections, selection functions, ionizing-field
uncertainty, and halo-to-halo scatter.

The central scientific target is not a census of ion labels. It is a
quantitative account of how mass, metals, momentum, and energy move between
galaxies and the cosmic web. Spectral lines are the measuring apparatus, and
their greatest power comes from combining transitions whose degeneracies are
different.
